\section{逆数による置換}\label{節:逆数対称性}

区間の端が $\frac1b$ と $b$ であるときは,\,$x=\frac1t$ でも両端を交換できる.ただし反転 $x=a+b-t$ と異なり,微分は
\[
\dd x=-\frac{1}{t^2}\dd t
\]
になる.関数の $x$ を逆数に変えるだけでは,正しい置換にならない.

特に積分が $\frac{\dd x}{x}$ を含むと,
\[
\frac{\dd x}{x}=-\frac{\dd t}{t}
\]
となり,同じ形に戻る.\,$b>1$ とすると,
\begin{equation}\label{式:逆数置換の対称性}
\int_{\frac1b}^b F(x)\frac{\dd x}{x}
 =\int_{\frac1b}^b F\left(\frac1x\right)\frac{\dd x}{x}
\end{equation}
である.したがって,\,$F(x)+F(\frac1x)=C$ なら
\[
2\int_{\frac1b}^b F(x)\frac{\dd x}{x}
 =C\int_{\frac1b}^b\frac{\dd x}{x}
 =2C\log b
\]
となる.

\begin{example}{逆数による計算}
次の定積分を求めよ.
\[
I=\int_{\frac13}^3\frac{1}{x(1+x^2)}\dd x.
\]
\end{example}

区間の両端が互いに逆数であり,分母に $x$ がある.そこで $F(x)=\frac{1}{1+x^2}$ と見て,\,$\frac{\dd x}{x}$ を残す.
\[
F\left(\frac1x\right)=\cfrac{1}{1+\cfrac{1}{x^2}}
 =\frac{x^2}{1+x^2},
\qquad F(x)+F\left(\frac1x\right)=1.
\]
式\eqref{式:逆数置換の対称性}を使えば
\[
2I=\int_{\frac13}^3\frac{\dd x}{x}
 =\left[\log x\right]_{\frac13}^3=2\log3,
\]
よって $I=\log3$ である.

この例を直接置換しても,
\[
I=\int_3^{\frac13}
\cfrac{1}{\cfrac1t(1+\cfrac1{t^2})}
\left(-\frac{1}{t^2}\right)\dd t
=\int_{\frac13}^3\frac{t}{1+t^2}\dd t
\]
となる.元の被積分関数との和は $\frac1x$ であり,上の計算と一致する.微分因子を忘れると,この和は現れない.

逆数の対称性を,いつもの左右の対称性として見直すこともできる.\,$x=\e^u$ と置くと,\,$\frac{\dd x}{x}=\dd u$,区間は $[-\log b,\log b]$ になる.さらに $x$ を $\frac1x$ に変える操作は,\,$u$ を $-u$ に変える操作になる.数直線上の普通の長さでは $\frac1b$ と $b$ は $1$ から等距離ではないが,対数で測った座標では左右対称なのである.

一般に,端が正で $ab=k$ なら,\,$x=\frac{k}{t}$ は $[a,b]$ の両端を交換する.この場合にも $\frac{\dd x}{x}=-\frac{\dd t}{t}$ が成り立つ.端を見て置換を考えるときは,和が一定か,積が一定かが手がかりになる.


